\documentclass{article}
\usepackage[T1]{fontenc}
\usepackage{graphicx,xcolor}
\title{An example manuscript for a Nature-inspired reading copy}
\author{First Author\textsuperscript{1}, Second Author\textsuperscript{2}\\
\small\textsuperscript{1}Department, University, City, Country\\
\small\textsuperscript{2}Institute, University, City, Country}
\date{}

\begin{document}
\maketitle
\begin{abstract}
This example demonstrates the input format accepted by the generator.
Replace the title, authors, affiliations and manuscript text with your own.
The generated layout is intended for reading and colleague review; it is
not an official journal template or a published article.
\end{abstract}

\section{Introduction}
Introduce the research problem, previous work and the contribution of the study.
Edit this source file rather than the generated reading copy.

\section{Results}
\subsection{Example finding}
Present the findings here. The figure below is a placeholder and requires
no external image file.

\begin{figure}[htbp]
\centering
\fbox{\rule{0pt}{35mm}\rule{0.85\textwidth}{0pt}}
\caption{An example figure placeholder. Replace this with your own artwork.}
\label{fig:example}
\end{figure}

\section{Discussion}
Discuss the findings, implications and limitations.

\section{Methods}
Describe the experimental and analytical methods.

\end{document}
